\section{Graded Exercise Bank}
\label{sec:ch3-exercises}

The following one hundred problems are cumulative.  The four tiers distinguish
technical fluency, theorem selection, synthesis, and open-ended mathematical
modelling.  Unless a problem states otherwise, assume sufficient smoothness for
the indicated derivatives and integrals to exist.

\subsection*{Tier I: Foundations (30 problems)}
\begin{enumerate}[label=\textbf{\arabic*.},leftmargin=*]
\item Sketch the vector field $\mathbf F=(x,y)$ and identify its source-like behavior.
\item Sketch $\mathbf F=(-y,x)$ and indicate the direction of circulation.
\item Compute $\nabla f$ for $f=x^2+3y^2-z$.
\item Compute $\nabla\cdot\mathbf F$ for $\mathbf F=(x^2,xy,z^2)$.
\item Compute $\nabla\times\mathbf F$ for $\mathbf F=(-y,x,0)$.
\item Compute $\nabla^2 f$ for $f=x^2+y^2+z^2$.
\item Evaluate $\int_C x\,ds$ on the line segment from $(0,0)$ to $(1,1)$.
\item Evaluate $\int_C (2x,2y)\cdot d\mathbf r$ from $(0,0)$ to $(1,2)$.
\item Parametrize the circle $x^2+y^2=a^2$ counterclockwise.
\item Find the unit tangent to $\mathbf r(t)=(t,t^2,0)$.
\item Find an oriented normal to the plane $x+2y+3z=6$.
\item Compute the flux of $\mathbf F=(0,0,1)$ through a horizontal disk of radius $R$.
\item Compute $\iiint_V 1\,dV$ for a rectangular box with sides $a,b,c$.
\item Write the cylindrical-coordinate volume element and explain the factor $r$.
\item Write the spherical-coordinate volume element and identify the polar-angle factor.
\item Determine whether $\mathbf F=(2x,2y)$ is conservative.
\item Find a potential for $\mathbf F=(y,x)$.
\item Evaluate the circulation of a gradient field around a closed curve.
\item State Green's theorem with positive boundary orientation.
\item State Stokes' theorem and its orientation rule.
\item State Gauss' theorem for a smooth vector field.
\item Compute the outward flux of $\mathbf F=(x,y,z)$ through the unit sphere using Gauss' theorem.
\item Compute the divergence of a constant vector field.
\item Compute the curl of a constant vector field.
\item Show directly that $\nabla\times(\nabla f)=0$ for $f=x^2y+z^3$.
\item Show directly that $\nabla\cdot(\nabla\times\mathbf A)=0$ for $\mathbf A=(yz,zx,xy)$.
\item Interpret the sign of $\nabla\cdot\mathbf v$ in a compressible fluid.
\item Interpret $\nabla\times\mathbf v=0$ in a simply connected fluid region.
\item Write the local conservation law for density $\rho$ and current $\mathbf J$.
\item Explain why the boundary of a boundary has no boundary in the integral theorems.
\end{enumerate}
\subsection*{Tier II: Intermediate synthesis (30 problems)}
\begin{enumerate}[label=\textbf{\arabic*.},leftmargin=*]
\item Find the flux of $\mathbf F=(x,y,0)$ across the outward boundary of the unit disk using Green's theorem.
\item Evaluate $\oint_C(-y\,dx+x\,dy)$ for the circle of radius $a$.
\item Determine whether $\mathbf F=(yz,xz,xy)$ is conservative and find a potential when possible.
\item Evaluate the work of $\mathbf F=(y,x+2y)$ along two different paths from $(0,0)$ to $(1,1)$ and compare.
\item Compute the surface area of the paraboloid $z=x^2+y^2$ above the disk $x^2+y^2\le 1$.
\item Compute the flux of $\mathbf F=(x,y,z)$ through the upper hemisphere and its base separately.
\item Use Gauss' theorem to find the flux of $\mathbf F=r^2\mathbf r$ through a sphere of radius $R$.
\item Use Stokes' theorem to compute the circulation of $\mathbf F=(-y,x,z)$ around a circle in the plane $z=c$.
\item Use Green's theorem to compute the area enclosed by an ellipse.
\item Derive the polar-coordinate area element from a Jacobian determinant.
\item Derive the cylindrical-coordinate scale factors.
\item Compute $\nabla^2(1/r)$ for $r\ne0$.
\item Verify that $u=e^{kx}\cos ky$ is harmonic in two dimensions.
\item Solve the steady one-dimensional diffusion equation with fixed endpoint temperatures.
\item For $\mathbf v=(ax,by,cz)$, determine the condition for incompressibility.
\item For rigid rotation $\mathbf v=\boldsymbol\Omega\times\mathbf r$, compute the vorticity.
\item Derive the integral continuity equation from the local one using Gauss' theorem.
\item Derive the local continuity equation from the integral form for arbitrary fixed volumes.
\item Show that a radial inverse-square field has constant flux through concentric spheres.
\item Find the electric field from $\phi=k/r$ away from the origin.
\item Compute the gravitational field produced by a spherically symmetric mass distribution outside its support.
\item Use integration by parts to derive Green's first identity.
\item Use Green's first identity to prove uniqueness for a Dirichlet problem under suitable assumptions.
\item Show that $\int_V |\nabla u|^2dV=-\int_V u\nabla^2u\,dV$ when $u=0$ on the boundary.
\item Explain why a nonzero circulation around a puncture does not contradict zero curl away from the puncture.
\item Compute the circulation of $\mathbf F=(-y/(x^2+y^2),x/(x^2+y^2))$ around a circle centered at the origin.
\item Find the divergence and curl of $\mathbf F=f(r)\mathbf r$.
\item Derive the radial Laplacian of a spherically symmetric scalar field.
\item Compute the mean value of a harmonic function over a circle for a simple polynomial example.
\item Verify numerically or symbolically one vector identity from the chapter.
\end{enumerate}
\subsection*{Tier III: Advanced analysis (20 problems)}
\begin{enumerate}[label=\textbf{\arabic*.},leftmargin=*]
\item Prove the path-independence theorem for conservative fields on a connected domain.
\item Give a counterexample showing that zero curl need not imply a global potential on a nonsimply connected domain.
\item Derive the divergence formula in orthogonal curvilinear coordinates from scale factors.
\item Derive the Laplacian in cylindrical coordinates.
\item Derive the Laplacian in spherical coordinates for a radial function.
\item Establish Green's second identity and state one consequence.
\item Use the divergence theorem to derive the electrostatic Poisson equation from Gauss' law.
\item Derive the wave equation from a first-order field system in a source-free region.
\item Prove that the flux of a curl through any closed smooth surface is zero.
\item Prove that the circulation of a gradient around any closed piecewise-smooth loop is zero.
\item Analyze the distributional divergence of $\mathbf r/r^3$ and determine its total source strength.
\item For a time-dependent control volume, state and derive the Reynolds transport theorem in a simple setting.
\item Derive the probability-current continuity equation from the Schrödinger equation.
\item Show how gauge transformation $\mathbf A\mapsto\mathbf A+\nabla\chi$ leaves $\mathbf B=\nabla\times\mathbf A$ unchanged.
\item Explain the differential-form pattern behind gradient, curl, divergence, and the identity $d^2=0$.
\item Prove a Helmholtz-type uniqueness statement under clearly stated decay and regularity assumptions.
\item Construct a divergence-free field from a vector potential and discuss gauge freedom.
\item Construct a curl-free field from a scalar potential and discuss topology.
\item Analyze boundary terms in an eigenvalue problem for the Laplacian.
\item Show that Laplacian eigenfunctions with distinct eigenvalues are orthogonal under suitable boundary conditions.
\end{enumerate}
\subsection*{Tier IV: Challenge and computational investigations (20 problems)}
\begin{enumerate}[label=\textbf{\arabic*.},leftmargin=*]
\item Derive the fundamental solution of the three-dimensional Laplacian and fix its normalization by flux.
\item Derive the two-dimensional logarithmic fundamental solution of the Laplacian.
\item Use separation of variables to solve Laplace's equation inside a disk with one nontrivial boundary Fourier mode.
\item Use separation of variables to solve the heat equation on an interval with homogeneous Dirichlet data.
\item Prove an energy-decay identity for the heat equation.
\item Derive Poynting's theorem from Maxwell's equations using a vector identity.
\item Relate circulation quantization around a puncture to phase winding of a complex wavefunction.
\item Analyze a magnetic vector potential for a uniform magnetic field in two different gauges.
\item Use Stokes' theorem to explain why gauge-related potentials yield the same magnetic flux.
\item Formulate Green, Stokes, and Gauss as instances of the generalized Stokes theorem.
\item Derive the weak form of Poisson's equation and identify the natural boundary term.
\item Prove a maximum principle for harmonic functions in a bounded domain at an appropriate level of rigor.
\item Derive the multipole expansion of a localized source through the first nontrivial terms.
\item Analyze the field and flux of a line singularity using cylindrical symmetry.
\item Explain how local differential laws can fail to determine global holonomy on a multiply connected domain.
\item Develop a computational experiment comparing direct surface integration with Gauss' theorem on a deformed closed surface.
\item Develop a computational experiment comparing direct line integration with Stokes' theorem on two spanning surfaces.
\item Prove conservation of total probability from the continuity equation under decaying or no-flux boundary conditions.
\item Show how the covariant divergence generalizes flux conservation in curvilinear coordinates.
\item Write a concise synthesis connecting vector calculus, gauge potentials, Berry phases, and later Quantum Hall theory.
\end{enumerate}

\subsection*{Selected hints and solution outlines}
\begin{description}[leftmargin=2.8cm,style=nextline]
\item[Foundation 22.] Since $\nabla\cdot\mathbf r=3$, the flux through the unit sphere is $3(4\pi/3)=4\pi$.
\item[Intermediate 8.] Choose the flat disk as the spanning surface; only the normal component of $\nabla\times\mathbf F$ contributes.
\item[Intermediate 17.] Integrate $\partial_t\rho+\nabla\cdot\mathbf J=0$ over $V$ and use Gauss' theorem.
\item[Intermediate 25.] The domain excludes the origin.  Topology, rather than a local derivative, controls the loop integral.
\item[Advanced 11.] Integrate over a ball with a small central ball removed, then take the inner radius to zero.
\item[Advanced 13.] Multiply the Schrödinger equation by $\psi^*$, subtract its conjugate, and collect the result as a divergence.
\item[Challenge 1.] Begin with spherical symmetry, solve away from the origin, and determine the coefficient from unit total flux.
\item[Challenge 6.] Dot Faraday's law with $\mathbf H$, dot Ampere--Maxwell with $\mathbf E$, subtract, and use $\nabla\cdot(\mathbf E\times\mathbf H)$.
\item[Challenge 8.] Compare symmetric gauge $\mathbf A=\tfrac12\mathbf B\times\mathbf r$ with a Landau gauge and exhibit the gauge function.
\item[Challenge 20.] Separate local curvature, encoded by field strength, from global phase information encoded by holonomy.
\end{description}
